\section{模型检验与敏感性分析}
\subsection{物理一致性与实际费用核验}
对问题一的144个区间以及问题二、三各策略的334天执行结果，逐段重算能量平衡、储电递推、功率与容量边界，并核对充放电互斥和跨日连续。所得能量及费用残差均低于$10^{-5}$的相应单位容差。对问题一，日初与日末储电量相同，还可将全天供给与消耗合并检验：
\begin{equation}
\sum_t(g_t+v_t)=\sum_t(\ell_t+w_t)
+\sum_t\left[(1-\eta)c_t+(\eta^{-1}-1)d_t\right].
\end{equation}
右侧最后一项为充放电损耗，不能因为购电量减少而忽略。

后两问的账单由实际计划、最终交付和紧急电量分别复算。未利用的已付费电量仍收费，取消电量不再供电，紧急电量只计一次5倍费用。零预测误差、光伏骤降、负载骤升以及满电、低电量等算例均满足物理约束。改变未来实测数据或尚未发布的预报，不改变此前计划和执行动作，表明回测遵守了决策时的信息限制。

\subsection{预计费用与实际费用的差异}
问题二主方案的日前情景预计紧急购电费合计为\QtwoExpectedEmergency\ 万元，而实际紧急购电费为\QtwoEmergencyCost\ 万元。预计值按每天前144段求和，不重复计入重叠的次日前瞻。两者存在明显差距，说明有限历史情景与情景内未来调度对真实缺电风险的估值仍偏乐观，不能把规划目标直接当作运行效果。

完全信息下界为\QtwoLower\ 万元，问题二主方案高于下界\QtwoGap\%。下界使用全部实测数据，并允许期末储电量在设备边界内自由选择，因此只能衡量改进空间。问题三的调整结算费用不低于按普通电价直接购买最终交付量的费用，故同一完全信息问题也为问题三提供有效下界。

\subsection{期末库存与参数敏感性}
不同策略期末储电量不同。为检查费用排序是否主要来自库存差异，采用低价时段补充储电的成本作为辅助估值：
\begin{equation}
\lambda=\frac{\min_t p_t}{\eta}=\QtwoInventoryPrice\ \text{元/kWh},
\qquad C^{\mathrm{adj}}=C-\lambda(E_{\mathrm{end}}-6000).
\end{equation}
此值用于比较，不改变实际账单。问题二主方案期末储电量为\QtwoEnd\ kWh，库存校正后相对单预测方案的费用差为\QtwoAdjustedDifference\ 元。问题三主方案期末为\QthreeEnd\ kWh，相对仅0点预报基准的校正后节省为\QthreeInventoryDifference\ 元，均未改变相应费用比较的方向。

问题一将单程效率由0.9改为$\sqrt{0.9}$，即往返效率由81\%提高至90\%，其他条件不变。重新优化后费用为\AltCost\ 元，相对主方案变化\AltDifference\ 元。这说明效率口径和能量损耗会影响成本，应在建模时明确。

问题二、三分别在四个指定日期采用14个情景，以及4800、7200 kWh的远端储电目标，从与主方案相同的历史信息和日初储电量出发重算。增加情景数量并非在每个日期都降低实际费用；远端目标变化对这些日期的费用影响很小。详细结果见附录\ref{app:sensitivity}，不据此推断全年效果或重新选择主参数。

对问题三另采用“取消部分原计划照付并追加0.5倍违约费”的结算解释，全部预报策略费用为\QthreeLiteralCost\ 万元。在允许免费舍弃余电时，保留多余供给即可避免新增违约费，因此此口径下减购没有经济优势，全年减购量为零。该对照说明结算规则会改变调整方向，不能把不同机制下的差异归结为预报质量。

\section{模型评价与改进方向}
\subsection{模型的优点}
首先，共同物理模型用能量守恒和储电递推连接所有时段，三问的差异通过信息条件和交易规则表达，避免为每个问题另设一套互不一致的调度机制。线性规划能够直接利用容量、功率与费用结构，不需要人为指定充放电时刻。

其次，日前计划与实际执行分开评价。计划购电、紧急补电及未利用供给都有明确去向和费用，跨日储电量连续，使比较反映连续运行的结果。无储能、单预测、预报组合及仅更新储能的对照，分别检验了设备配置、误差处理和交易调整的作用。

最后，所得策略具有可解释性：富余光伏为储能提供能量，电价差决定跨时段转移是否划算，紧急补电的高价格推动日前保留余量，而调整结算又限制了频繁改动计划的收益。

\subsection{局限与改进方向}
历史预测方法较简单，有限情景可能遗漏季节变化和极端天气下的误差，日内权重更新也只能调整已有情景的可信程度。问题二预计紧急费用显著低于实际费用，提示后续改进应优先检查误差覆盖与风险估值，而不能仅通过增加求解复杂度判断方法更好。

模型未计电池衰减、自放电和设备动态，小时预报的线性插值也不能恢复更快的光伏波动。若获得相应数据，可在现有模型中增加循环损耗成本、测量延迟或更细的功率约束，并重新比较实际费用。

规划期内各情景的未来动作属于前瞻近似，当前动作才满足同一信息下的一致性要求。因此本文没有证明完整多阶段随机控制的全局最优性。结论来自所给单年数据、指定设备和结算假设；增加其他发布时间、推广到其他年份或改变设备规模，都需要对应数据重新验证。

\section{结论}
在已知全天输入时，问题一采用144段线性规划。最优购电量为\QoneGrid\ kWh、费用\QoneCost\ 元，较无储能基准降低\SavingPercent\%。这一结果来自富余光伏消纳与跨时段购电安排，同时满足首尾储电量相等及设备边界。

在存在预测误差时，多情景日前计划与滚动储能控制的334天实际费用为\QtwoCost\ 万元，较无储能多情景购电降低\QtwoNoStorageSaving\%，紧急购电发生于\QtwoDays\ 天。历史误差情景改善了本次对照中的购电决策，但实际紧急费用仍高于规划估值，误差风险不能由单一预测或规划目标代替评价。

在允许日内修订时，6、12、18点预报全部采用的方案费用为\QthreeCost\ 万元，较仅0点预报降低\QthreeSaving\%。八种组合的比较支持在本次条件下采用这三个发布时间，同时“仅更新预报、不修改购电”的策略费用略增，说明新信息的价值取决于能够据此作出的决策。储能调度、信息更新与交易调整应在统一的实际费用口径下共同评价。
